% Shared preamble for the CIG-AMF formal-research papers.
% IEEEtran conference class; compiles on Overleaf without modification.
\documentclass[conference]{IEEEtran}
\IEEEoverridecommandlockouts

\usepackage{cite}
\usepackage{amsmath,amssymb,amsfonts}
\usepackage{amsthm}
\usepackage{algorithm}
\usepackage{algorithmic}
\usepackage{graphicx}
\usepackage{textcomp}
\usepackage{xcolor}
\usepackage{booktabs}
\usepackage{url}
\usepackage[hidelinks]{hyperref}

\def\BibTeX{{\rm B\kern-.05em{\sc i\kern-.025em b}\kern-.08em T\kern-.1667em\lower.7ex\hbox{E}\kern-.125emX}}

\newtheorem{theorem}{Theorem}
\newtheorem{lemma}[theorem]{Lemma}
\newtheorem{proposition}[theorem]{Proposition}
\newtheorem{corollary}[theorem]{Corollary}
\theoremstyle{definition}
\newtheorem{definition}[theorem]{Definition}
\newtheorem{remark}[theorem]{Remark}
\newtheorem{example}[theorem]{Example}

% A machine-checked statement carries its Lean declaration name.
\newcommand{\lean}[1]{\textsf{\footnotesize #1}}
\newcommand{\leanbox}[1]{%
  \par\vspace{2pt}\noindent
  {\footnotesize\textsf{Lean: #1}}\par\vspace{2pt}}

% Evidence classes, used consistently and never interchangeably.
\newcommand{\PROVED}{\textsc{proved}}
\newcommand{\CERT}{\textsc{certificate}}
\newcommand{\NUM}{\textsc{numerical}}
\graphicspath{{../_figures/}}


\begin{document}

\title{Supplement to ``The Price of a Frozen Surrogate in Support-Aware Compression'':\\
Proofs, Formal Declarations and Conditional Results}

\author{\IEEEauthorblockN{Vu Phan An Nguyen}
\IEEEauthorblockA{\textit{Rider University} \\
Lawrenceville, NJ, USA}}

\maketitle

\begin{abstract}
This supplement gives, for each theorem of the main paper, the proof architecture and the
Lean~4 declarations it rests on. It also collects the material the main text leaves out, each
with its exact status: the conditional rate of the implemented learner, the confidence-sequence
machinery, two heteroscedastic families, the repair criterion, and the older constructions that
the centred theorem supersedes. Sections S1--S9 follow the main text, and S10 records the build; theorems are named as
there.
\end{abstract}

\section*{Conventions}

All declarations live in \texttt{work/p13\_conjectures}. A module is accepted when it compiles
with exit status $0$, contains no \texttt{sorry} and no \texttt{native\_decide}, and every audited
axiom closure lies within \{\texttt{propext}, \texttt{Classical.choice}, \texttt{Quot.sound}\}.
Acceptance is checked by a frozen gate that copies every module and its local imports into a
fresh directory and builds from an empty cache (S10). A theorem that takes a hypothesis is
reported with that hypothesis named. It is not called a result about an algorithm until the
hypothesis is discharged at that algorithm.

\section{S1. The model, the decision and its alternatives}

\paragraph{The loss is half a range} The coefficient vector $a_{T}$ is defined from $F$, the
reference and $T$. The identity making it the residual is proved, as is the loss formula, and
the corpus's \lean{selectedRadius} is identified with half the oscillation of that residual.

\leanbox{\lean{residual\_eq\_inner}, \lean{loss\_eq\_half\_range},
\lean{selectedRadius\_eq\_compressionLoss} in \lean{LeanPaperAModelBridgeV1}}

\paragraph{The class and the marginals} The centred class is exactly where both
uniform-product marginals recover their coordinates. On it, the residual after retaining $A$
is $\lambda_{B}$, so the loss is half of arm $B$'s range. The unsafe condition carries its
tolerance and is strict.

\leanbox{\lean{centred\_iff\_marginals}, \lean{residA\_eq}, \lean{rangeB\_eq\_osc},
\lean{lossRetainA\_eq}, \lean{unsafeA\_iff\_range\_gap}, \lean{unsafeA\_zero\_iff},
\lean{unsafe\_subset\_closed} in \lean{LeanPaperACentredModelV1}}

\paragraph{The alternative set, exactly} Retaining $A$ is $\varepsilon$-unsafe exactly when some
pair $(p,q)$ of arm-$B$ cells satisfies two simultaneous strict linear inequalities. Forward,
take the cells realising arm $B$'s extrema. Backward, any pair's gap is at most the range. The
two inequalities are simultaneous, from $|x|<y\iff x<y\wedge -x<y$. Covering the set with the
two as \emph{alternatives} contains the true world itself, and a certificate built on that
cover collapses to zero.

\leanbox{\lean{range\_gap\_iff\_exists\_pair}, \lean{pair\_region\_linear},
\lean{unsafe\_iff\_union\_of\_polyhedra}, \lean{polyhedron\_nonvacuous} in
\lean{LeanPaperARegionV1}}

\paragraph{General regions, and the closure} For a general retained set and competitor, pinning
the four extremum indices cuts out a polyhedron on which the doubled loss difference is one
linear form. The alternative set is the union of those regions, in both directions. On a
region that contains a strictly bad point, the infimum over the strict part equals the infimum
over the closure. The proof moves along the segment towards that point, where the energy moves
by at most $t$ times an explicit constant and the segment is strictly bad for every $t>0$. A
region with an empty strict part is dropped, because closing it would turn a tie into an
alternative. A coordinate with zero weight leaves the objective, so the inner problem is
unconstrained in it.

\leanbox{\lean{loss\_eq\_of\_extrema}, \lean{region\_form\_is\_linear}, \lean{alt\_iff\_exists\_region},
\lean{exists\_extrema}, \lean{strict\_inf\_eq\_closed\_inf}, \lean{segment\_stays\_strict},
\lean{energy\_segment\_bound}, \lean{strict\_empty\_has\_no\_alternatives} in
\lean{LeanPaperAAlternativesV1}; \lean{exists\_strict\_near}, \lean{lower\_bound\_transfers},
\lean{empty\_strict\_contributes\_nothing}, \lean{objective\_indep\_of\_zero\_weight},
\lean{zero\_weight\_unbounded\_below\_direction} in \lean{LeanPaperAOracleRegionV1}}

\paragraph{Certificates at two levels} A KKT point with non-negative multipliers minimises the
quadratic over its region against every feasible competitor. The Hessian is \emph{not} assumed
invertible, since a design that leaves a direction unmeasured has a singular one. A prior on
alternatives that satisfies the dual inequality at every admissible query bounds every design,
so ``optimal'' does not mean ``optimal within a searched family''. A certificate whose
condition implies safety is never cheaper than the decision.

\leanbox{\lean{kkt\_region\_minimum}, \lean{quad\_expand}, \lean{design\_upper\_bound\_of\_dual},
\lean{design\_value\_eq\_of\_primal\_dual}, \lean{certification\_costs\_at\_least\_as\_much} in
\lean{LeanPaperAOracleCertificateV1}}

\paragraph{The divergence is the energy} For every centred candidate, the Gaussian divergence
of each query's observation law is $\frac{(\theta_{e}-\lambda_{e})^{2}}{2\nu_{e}}$, so the
energy of a design is the design-weighted divergence. The game over the strict alternatives is
the game of the observation law.

\leanbox{\lean{energy\_eq\_sum\_gaussKL}, \lean{gaussKL\_eq\_integral} in \lean{LeanPaperAStrictV1};
\lean{gaussianKLIdentification\_holds} in \lean{LeanPaperAGaussianKLV1};
\lean{actualKLValue\_eq} in \lean{LeanPaperAEndgameV1}}

\section{S2. Separation (Theorem ``Separation'')}

\paragraph{The frozen value} Shifting the four frozen cells by $(+\frac g4,-\frac g4,-\frac g4,
+\frac g4)$ closes the contrast, and this witness is centred. So
$J_{\rm fr}(w)\le\frac{g^{2}}{32}a(w)$. The uniform-anchor design attains $\frac{g^{2}}{32}$ by
Cauchy--Schwarz, and $g=\frac34$ gives $\frac{9}{512}$. The restriction to $\Theta_{0}$ raises
the infimum, so the upper bound needs the witness to be in the class, which it is.

\leanbox{\lean{thetaC\_centred}, \lean{frozen\_margin}, \lean{energy\_lamStar}, \lean{lamStar\_frozenC₀},
\lean{frozenValueC₀\_le}, \lean{frozenGameC₀\_eq}, \lean{frozenValueC\_le\_restricted} in
\lean{LeanPaperACentredV1}}

\paragraph{The true value} Upper: some pair carries at most $\frac1r$ of any allocation, and the
pair shift by $\frac12$ is a centred alternative of energy $\frac{\mathrm{pairWeight}}{8M}$. For
every $\delta>0$ the shift by $\frac12+\delta$ is strictly alternative. Lower: at
$w_{e}=\nu_{e}/(4+2rM)$ the energy is the squared distance over $2(4+2rM)$, and a reversal
forces the four-cell contrast to move by $\frac34$.

\leanbox{\lean{exists\_light\_pair\_le}, \lean{trueGameC₀\_le}, \lean{pairShiftC\_strict\_alt},
\lean{wCert\_isAlloc}, \lean{energy\_wCert}, \lean{contrast\_movesC}, \lean{trueValueC₀\_wCert\_ge},
\lean{trueGameC₀\_ge}, \lean{centred\_separation} in \lean{LeanPaperACentredV1};
\lean{trueValueC₀\_le\_trueValueStrict}, \lean{trueValueStrict\_le\_shift} in
\lean{LeanPaperAStrictV1}; \lean{strict\_design\_upper}, \lean{strict\_game\_bounds} in
\lean{LeanPaperAEndgameV1}}

\NOTEPAR{Earlier constructions. The heteroscedastic families of \lean{LeanPaperASeparationV1}
separate only as $M$ grows. They are cited in the main text for the underpricing direction, which
the centred construction does not show. The $r$-pair games of \lean{LeanPaperARPairV1},
\lean{LeanPaperARGamesV1} and \lean{LeanPaperAFrozenV1} bracket the ratio between
$\frac{1}{128(4+2rM)}$ and $\frac{16}{Mr}$ on an uncentred world, against a restricted-coordinate
surrogate and a frozen one. The centred theorem gives the sharper bracket on one class.}

\leanbox{\lean{no\_constant\_factor\_on\_families} in \lean{LeanPaperASeparationV1};
\lean{rpair\_ratio\_le}, \lean{rpair\_ratio\_ge} in \lean{LeanPaperARPairV1};
\lean{rpair\_games\_theta} in \lean{LeanPaperARGamesV1}; \lean{frozen\_games\_theta} in
\lean{LeanPaperAFrozenV1}}

\section{S3. Starvation (Theorem ``Starvation'')}

$J_{\rm fr}(w)\le\frac{9}{512}a(w)$, and $a(w)$ plus the middle mass is $1$. At the optimum
$a(w)=1$. Within $\varepsilon$ of it, the middle mass is at most $\frac{512}{9}\varepsilon$.
Some pair then carries at most $\frac1r$ of that, and the strict pair shift bounds the true
value by $\frac{\mathrm{pairWeight}}{8M}$, giving $\frac{64\varepsilon}{9rM}$. The strict version
of the zero statement uses the shift by $\frac12+\delta$ at every $\delta>0$.

\leanbox{\lean{aMass\_add\_middle}, \lean{starvation\_at\_optimum₀}, \lean{near\_optimal\_starvation₀},
\lean{frozen\_optimum\_trueValueC₀\_zero}, \lean{near\_optimal\_trueValueC₀\_le} in
\lean{LeanPaperACentredV1}; \lean{starved\_trueValueStrict\_zero},
\lean{frozen\_optimum\_trueValueStrict\_zero} in \lean{LeanPaperAStrictV1}}

\NOTEPAR{The homogeneity route. Before the centred value was exact, starvation was derived from
positive homogeneity of the surrogate in the anchor weights, giving the middle mass bound
$\varepsilon/(\Gamma_{\rm fr}-\varepsilon)$ with no monotonicity modulus. It is compiled and
superseded.}

\leanbox{\lean{middle\_mass\_le\_of\_homogeneous}, \lean{middle\_mass\_le\_of\_near\_optimum} in
\lean{LeanPaperAHomogV1}}

\section{S4. The price of repair (Theorems ``The repair frontier'', ``The sacrifice law'',
``The forcing law'')}

\paragraph{The frontier} The energy is linear in the design, so each value is an infimum of
linear functions and is concave. Mixing the uniform-anchor design with the certifying design
keeps $(1-\beta)$ of the surrogate and gains $\beta$ of the true lower bound. The upper side is
starvation with $\varepsilon=\beta\Gamma_{\rm fr}$. On the strict game the pointwise exchange
$J_{\rm true}(w)\le\mathrm{sac}(w)/(8rM)$ comes from the strict pair shift at every margin
$\frac12+d$, followed by $d\to0$.

\leanbox{\lean{mixW\_isAlloc}, \lean{energy\_mixW}, \lean{value\_concave}, \lean{frontier\_upper},
\lean{frontier\_lower}, \lean{frontier\_two\_sided}, \lean{frontier\_at\_zero} in
\lean{LeanPaperAFrontierV1}; \lean{limit\_shift}, \lean{exchange\_strict}, \lean{strict\_frontier}
in \lean{LeanPaperASacrificeV1}}

\paragraph{The exact slope} Let $v_\beta=(1-\beta)w_{\rm unif}+\beta u$, where $u$ is uniform on
the $2r$ middle cells. Its frozen value is at least $(1-\beta)\frac{9}{512}$, by concavity, so
$\mathrm{sac}(v_\beta)\le\beta$. Its weights are $\frac{1-\beta}{4}$ on each frozen cell and
$\frac{\beta}{2r}$ on each middle cell, of variance $M$. A strict alternative reverses the order of
the ranges. Split by which two arm-$B$ cells realise its extrema: two middle cells (the contrast
to close is $1$), one anchor and one middle cell ($\frac78$), both anchors ($\frac34$), or a
single cell, when the arm-$A$ cells alone must move by $1$. In each case the contrast is spread
over at most four moved cells, and the weighted Cauchy--Schwarz inequality
$g^{2}/\sum_i c_i^{-1}\le\sum_ic_ix_i^{2}$, with $c_i$ the weight over twice the variance, bounds
the energy below. With $X=rM$ and $0<\beta\le\frac12$ every case is at least
$\beta/(8X+128\beta)$. The upper side is the pointwise exchange. The two ends differ by
$\frac{128\beta^{2}}{8X(8X+128\beta)}\le\frac{2\beta^{2}}{X^{2}}$, so the slope at zero is
exactly $\frac1{8rM}$. On the strict game the lower side transfers, since restricting to strict
alternatives raises the infimum.

\leanbox{\lean{uMid}, \lean{uMid\_isAlloc}, \lean{vB}, \lean{vB\_isAlloc}, \lean{vB\_frozen},
\lean{sq\_tangent}, \lean{weighted\_four}, \lean{weighted\_two}, \lean{bcell\_cases},
\lean{trueValue\_vB\_ge}, \lean{frontier\_slope}, \lean{slope\_gap} in
\lean{LeanPaperARepairSlopeV1}; \lean{trueValueC₀\_le\_trueValueStrict} in \lean{LeanPaperAStrictV1}}

\paragraph{The sacrifice law} Fix pair $j$ and $d>0$. The world that shifts pair $j$ by
$\frac12+d$ is a strict centred alternative. Only the two cells of pair $j$ separate it from
$\theta_{C}$, each with divergence $\frac{(1/2+d)^{2}}{2M}$. The change of measure for any
randomised adaptive policy gives
$\sum_{q}\mathbb{E}[N_{q}]\,\mathrm{KL}_{q}\ge(1-\delta)\log\frac1\delta-\log2=D$. The dummy
query carries nothing, and the real terms are the energy of the expected counts. So
$D\le N_{\rm pair}\frac{(1/2+d)^{2}}{2M}$ for every $d>0$, and $N_{\rm pair}\ge8MD$. Summing over
the $r$ pairs gives $8rMD$ middle queries. The surrogate's value is at most $\Gamma_{\rm fr}$
times the non-middle mass, so the middle mass of an allocation is at most its relative
sacrifice. Applied to the average design $\bar w$, this gives $N\,\mathrm{sac}(\bar w)\ge8rMD$.
The law is stated for $D>0$, which also makes $N$ positive and $\bar w$ an allocation; every
$\delta<\frac18$ gives $D>0$ (\lean{hD\_of\_small} in \lean{LeanPaperARepairGaussV1}).

\leanbox{\lean{demand}, \lean{sacrifice}, \lean{sacrifice\_eq\_relative},
\lean{sacrifice\_le\_one}, \lean{middle\_le\_sacrifice}, \lean{expCnt}, \lean{demand\_le\_shift},
\lean{pair\_demand}, \lean{middle\_demand}, \lean{sacrifice\_cost\_law} in
\lean{LeanPaperASacrificeV1}; \lean{change\_of\_measure}, \lean{transcript\_kl\_le\_qsum} in
\lean{LeanPaperALowerKLV1}}

\paragraph{The upper constant, and satisfiability} The strict game is at least
$\frac{9}{128(4+2rM)}$, so $1/\Gamma_{\rm true}\le\frac{128(4+2rM)}{9}$. Its ratio to the lower
constant $8rM$ is $\frac{32}{9}+\frac{64}{9rM}$. The ideal learner is truncated at a horizon
chosen by Markov's inequality, so that it has stopped with probability at least $1-\delta$. Its
answer event then has probability at least $1-2\delta$ at the truth and at most $2\delta$ at
every strict alternative. It therefore satisfies every hypothesis of the law at level
$2\delta$, and its expected number of real queries is at most its expected stopping time.

\leanbox{\lean{upper\_constant}, \lean{constant\_gap} in \lean{LeanPaperASacrificeV1};
\lean{ideal\_sacrifice}, \lean{sacrifice\_hypotheses\_satisfiable} in
\lean{LeanPaperASacrificeIdealV1}; \lean{markov\_horizon} in \lean{LeanPaperA7IdealFamilyV1}}

\paragraph{Unbounded stopping (Theorem ``The law for unbounded stopping'')} The stopping policy
issues the free dummy after its stop, so one kernel policy serves every horizon. Its output read
at the first stopping time is \emph{false} (``not $A$'') on a run that has not stopped. Hence at a
strict alternative the event ``answers $A$ by horizon $n$'' lies inside ``stopped and answered
$A$'', which has probability at most $\delta$, and no stopping at the alternatives is needed. At
$\theta_{C}$, a horizon where the unstopped mass is at most $\delta$ leaves the wrong-or-unstopped
mass at most $2\delta$. The truncation therefore meets the finite-horizon law at level $2\delta$.
Expected counts only grow with the horizon, so the unbounded run's counts dominate the
truncation's middle demand, and the middle mass of any allocation is at most its sacrifice.

\leanbox{\lean{ansT}, \lean{outT\_unfired}, \lean{CorrectA}, \lean{trunc\_A}, \lean{expCntSt},
\lean{expCnt\_le\_St}, \lean{sacrifice\_law\_stopped}, \lean{correctA\_ofFinite} in
\lean{LeanPaperASacrificeStoppedV1}; \lean{stoppedPol}, \lean{outT\_first}, \lean{ofFinite} in
\lean{LeanPaperBStoppedLawV1}}

\paragraph{Spreading the budget (Theorem ``Spreading the budget'')} The padded policy runs
$\pi_{0}$'s kernels for $n$ rounds and then a deterministic round-robin over the four frozen cells (two $A$ cells and two $B$ anchors). Its first $n$ rounds therefore have the law of $\pi_{0}$: the two policies share their
kernels before round $n$, and the $n$-round marginal of the longer history law is the $n$-round
law. The answer event read on the first $n$ rounds keeps its probabilities. The expected count of
a cell grows at a padded round by the mass that round's kernel puts on it, which is one on the
scheduled $A$ cell. After $4k$ rounds every $A$ cell has gained exactly $k$. The padded average
design is then the mixture $(1-\beta)\,w_{\rm unif}+\beta\,\bar w_{0}$ with
$\beta=N_{0}/(N_{0}+4k)$. Concavity of the frozen value and $J_{\rm fr}(w_{\rm unif})\ge\frac{9}{512}$
give $\mathrm{sac}\le\beta\,\mathrm{sac}(\bar w_{0})$.

\leanbox{\lean{sacrifice\_mix\_le}, \lean{padQ}, \lean{padPol}, \lean{padPol\_lt}, \lean{padPol\_ge},
\lean{pad\_hist\_n}, \lean{pad\_real}, \lean{padCnt}, \lean{lintegral\_cnt\_pad}, \lean{padCnt\_A},
\lean{padCnt\_mid}, \lean{expCnt\_pad\_A}, \lean{expCnt\_pad\_mid}, \lean{dilution},
\lean{dilution\_small} in \lean{LeanPaperASacrificeDiluteV1}; \lean{hist\_congr},
\lean{hist\_map\_restr} in \lean{LeanPaperBStoppedLawV1}; \lean{lintegral\_qsum\_succ} in
\lean{LeanPaperALowerKLV1}}

\paragraph{The forcing law} The counts are bounded by the horizon, so $T$ and $\varphi(T)$ are
bounded and measurable. Jensen for the concave $\varphi$ gives
$\mathbb{E}\varphi(T)\le\varphi(\mathbb{E}T)$, and the middle demand closes it. For
$\varphi(t)=ct^{a}$, concavity of $t^{a}$ on $[0,\infty)$ and monotonicity of $x\mapsto x^{1/a}$
invert it.

\leanbox{\lean{totReal}, \lean{midReal}, \lean{integral\_midReal}, \lean{integrable\_totReal},
\lean{forcing\_law}, \lean{power\_forcing\_law} in \lean{LeanPaperASacrificeV1}}

\NOTEPAR{Superseded: the cumulative-exploration shape. Before the forcing law, the same shape
was assembled from four links, the per-round escape probability and the middle mass of an
$\varepsilon$-optimal design among them, with the change-of-measure demand taken as an input.
It is compiled, takes that input, and is no longer needed.}

\leanbox{\lean{inv\_sq\_shift\_sum\_le}, \lean{expected\_bad\_rounds\_bounded},
\lean{cumulative\_exploration\_bound} in \lean{LeanPaperA2ChainV1};
\lean{surrogate\_cost\_quadratic}, \lean{surrogate\_omega\_log\_squared} in
\lean{LeanPaperASurrogateV1}}

\section{S5. Learners (Theorems ``The ideal learner'' and ``What following the surrogate
costs'')}

\paragraph{The ideal learner} Its oracle maximises the exact value $J$ over a refining net, and
its stop uses the exact infimum of the ridge statistic
$\frac12\sum_e N_e^2(\widehat\theta_e-\lambda_e)^2/((a+N_e)\nu_e)$,
so its stopping-certificate gap is zero by construction. This is \lean{glr} in
\lean{LeanPaperA7GLRMeasV1}; despite that identifier, it is not the ordinary Gaussian
log-GLR with coefficient $N_e/\nu_e$. The main text gives the full query and stopping rule.
\begin{itemize}
\item \emph{Correctness.} Ville's inequality applies at the learner's own statistic, so the
probability of stopping with a wrong answer is at most $(|\mathrm{Ans}|-1)\delta$.
\item \emph{Upper bound.} On the rounds where the confidence event holds, the statistic
accumulates at $\Gamma-\eta$ for every $\eta>0$, with no box on the worlds. The event fails at
round $s$ with probability at most $s^{-4}$. Summing these gives, for every
$\eta\in(0,\Gamma)$, some $A_{\eta}$ with
$\mathbb{E}\tau_{\delta}\le\log(1/\delta)/(\Gamma-\eta)+A_{\eta}$ at every $\delta\in(0,1)$.
\item \emph{Lower bound.} The stopped learner, which queries a fresh dummy after stopping, is
coupled path by path with the unstopped run. Its real-cell count within any horizon is at most
$\tau_{\delta}$, and the policy lower bound applies to it at level $2\delta$.
\item \emph{Identification.} On the family, the learner's game is the Gaussian-divergence game
over the strict alternatives.
\end{itemize}

\leanbox{\lean{ideal\_expect\_bound}, \lean{ideal\_upper\_rate}, \lean{ideal\_finite},
\lean{ideal\_correct} in \lean{LeanPaperA7IdealUpperV1}; \lean{ideal\_direct} in
\lean{LeanPaperA7IdealDirect2V1}; \lean{runS\_eq}, \lean{cnt\_le\_tau}, \lean{ansEv\_alt\_le},
\lean{ansEv\_true\_ge} in \lean{LeanPaperA7IdealStoppedV1}; \lean{GamT\_fam},
\lean{ideal\_family\_lower}, \lean{ideal\_family\_lower\_rate}, \lean{ideal\_family\_finish},
\lean{ideal\_family\_correct} in \lean{LeanPaperA7IdealFamilyV1}; \lean{family\_policy\_lower},
\lean{actualKLValue\_le\_game} in \lean{LeanPaperAFamilyLowerV1}}

\paragraph{Following the surrogate, at the truth} Tracking a fixed design that is zero on the
middle cells, the pair's count up to $\min(n,\tau_{\delta})$ is at most
$4+2(\tau_{\delta}/s+s)/K$ along every path, for every $s>0$. The stopped run's counts equal the
unstopped run's at $\min(n,\tau_{\delta})$. The stopped change of measure at level $2\delta$
gives $\mathbb{E}N_{\rm pair}/(2M)\ge(1-2\delta)\log(1/(2\delta))-\log2$. Taking
$s=\sqrt{\mathbb{E}\tau_{\delta}}+1$ gives the square.

\leanbox{\lean{klq\_pair}, \lean{starved\_lower\_at}, \lean{starved\_cost\_log\_sq},
\lean{surrogate\_optimum\_cost}, \lean{surrogate\_optimum\_exists}, \lean{uniform\_surrogate\_cost},
\lean{surrogate\_vs\_ideal} in \lean{LeanPaperA6SurrogateCostV1}; \lean{fCount\_le\_amgm},
\lean{fCount\_le\_sqrt} in \lean{LeanPaperA6FixedDesignV1}; \lean{cnt\_eq\_at\_min},
\lean{ansEv\_alt\_leG}, \lean{ansEv\_true\_geG} in \lean{LeanPaperA7StoppedGenV1}}

\paragraph{Following the surrogate, at the estimate} Consider an estimate whose extrema are the
four anchors and which retains $A$. There the unrestricted surrogate's value is
$\frac{g(\hat\mu)^{2}}{32}$, and every exact optimum is zero on the middles. Forcing gives
every count order $\sqrt t$. So after a deterministic time $s_{0}$, the learner's confidence
event places every estimate within $\frac1{16}$ of the truth, which makes the anchors the
extrema. The event fails at round $s$ with probability at most $s^{-4}$. The expected number of
misleading rounds is therefore at most $s_{0}+2$ at every horizon, and it enters the pair's count
additively.

\leanbox{\lean{est\_cost\_log\_sq}, \lean{surrogate\_learner\_cost}, \lean{exact\_oracle\_cost},
\lean{surrogate\_learner\_vs\_ideal}, \lean{est\_lower\_at} in \lean{LeanPaperA6EstCostV1};
\lean{frozenGameAt\_eq}, \lean{starvation\_at}, \lean{surD\_exact}, \lean{optimal\_oracle\_middle\_zero},
\lean{close\_extreme} in \lean{LeanPaperA6SurrogateAtV1}; \lean{radius\_eventually},
\lean{close\_of\_good}, \lean{badAll\_expect} in \lean{LeanPaperA6BadRoundsV1};
\lean{oCount\_le\_amgm}, \lean{oCount\_ge} in \lean{LeanPaperA6OracleLearnerV1}}

\paragraph{Two surrogates} The centred surrogate at the truth and the unrestricted surrogate at
an estimate agree at $\theta_{C}$, where both are $\frac{9}{512}$. Now shift both arm-$A$ cells by
$u$. The unrestricted value of the uniform-anchor design stays $\frac{9}{512}$, because the
frozen margin ignores a common shift. The centred value must restore $x_{0}+x_{1}=0$. Write
$d_{i}$ for the displacements. Then $d_{0}+d_{1}=-2u$ and $(d_{0}-d_{1})-d_{2}+d_{3}\ge\frac34$.
Cauchy--Schwarz with weights $(\frac12,1,1)$ gives $\sum d_{i}^{2}\ge2u^{2}+\frac{9}{64}$, hence
a value of at least $\frac{9}{512}+\frac{u^{2}}{4}$.

\leanbox{\lean{frozenCentredAt}, \lean{muShift}, \lean{frozenGameAt\_thetaC},
\lean{frozenAt\_unif\_shift}, \lean{energy\_unif}, \lean{centred\_unif\_shift\_ge},
\lean{two\_surrogates\_differ} in \lean{LeanPaperA6TwoSurrogatesV1}}

\section{S6. The implemented learner, and what its rate still takes}

\paragraph{The algorithm} The shipped learner, \texttt{EXP\_paperA7\_learner.py}, sees only
counts, observations and known variances. At each refresh it brackets each answer's value by
cutting planes at the estimate. It keeps the answers whose upper bracket survives
$\max_{S}l_{S}-2Lr_{t}$, tracks the selected answer's design with forced exploration, and stops
when a certified lower bound on every region's infimum exceeds the threshold. The certificate is
weak duality, valid whatever the optimiser does. A primal objective is an upper bound and is
never used to stop. The cutting-plane tolerance tightens as $1/n^{2}$ down to a floor of
$10^{-6}$, with an 80-iteration cap, and weights are floored at $10^{-9}$. The alternatives are
those of the box model $[-8,8]^E$, so the certificate is valid for worlds in the box. The true
world is not in the scope of \texttt{run}, and the
self-check inspects the signature.

\paragraph{The conditional chain} For this learner the stopping chain is a composition of proved
links whose inputs are hypotheses about abstract processes.
\begin{itemize}
\item On a window where the confidence event holds, the statistic accumulates at
$\Gamma^{*}$ minus the window's radius, selection slack, tracking discrepancy and
initialisation loss.
\item Not having stopped forces a confidence failure in the window. With the window
$\lceil\sqrt t\rceil$ and failure mass $s^{-4}$, $\Pr(\tau>t)\le8t^{-3/2}$ and
$\Pr(\tau=\infty)=0$.
\item The expectation bound $\mathbb{E}\tau\le\log(1/\alpha)/\Gamma^{*}+19$ is proved
\emph{under the hypothesis} that each truncated expectation is at most
$t_{0}+\sum_{t\ge t_{0}}\Pr(\tau>t)$ (the input \texttt{hEsum}). It is not an unconditional
statement about the learner. The joined chain pays the window head,
$2\Gamma\sqrt t\le\varepsilon Gt+\Gamma^{2}/(\varepsilon G)$. So it lands at
$1/((1-\varepsilon)\Gamma)$ and does not reach $1/\Gamma$.
\end{itemize}
Three inputs remain for the shipped learner: the $s^{-4}$ confidence bound at its own
statistic; the accumulation hypotheses for its actual statistic, including the
stopping-certificate gap; and the passage to the $\mathbb{N}_{\infty}$ hitting time. For the
ideal learner all three are discharged (S5).

\leanbox{\lean{good\_window\_rate}, \lean{stopping\_subset\_window\_failure},
\lean{window\_union\_bound}, \lean{sqrt\_window\_tail}, \lean{tail\_forces\_finite\_stopping},
\lean{sqrt\_tail\_sum\_le}, \lean{expected\_stopping\_bound\_at\_level}, \lean{a7\_rate\_limsup} in
\lean{LeanPaperA7TailV1}; \lean{sqrt\_head\_split}, \lean{chain\_expected\_bound},
\lean{a7\_chain\_rate}, \lean{chain\_rate\_is\_deficient} in \lean{LeanPaperA7ChainV1};
\lean{increment\_deficit\_gives\_eta}, \lean{expected\_stopping\_bound\_general},
\lean{rate\_limsup\_eta}, \lean{eta\_to\_zero} in \lean{LeanPaperA7EtaV1};
\lean{sInf\_Qstat\_ge}, \lean{perturbationCost\_div\_le}; \lean{expectation\_le\_split} in
\lean{LeanPaperClosuresV1}}

\paragraph{The bracket width} As shipped, the oracle's bracket settled at $8.0\times10^{-3}$
against $\Gamma_{\rm comp}\approx2.3\times10^{-2}$. Taking the running minimum of the relaxation
values, damping the step and solving the regional dual to tolerance reduced it to
$5.0\times10^{-7}$, flat in both the cutting-plane count and the dual budget. A bracket of width
$\varepsilon$ gives $0\le\Gamma-J(w)\le\varepsilon$, and so a certified rate of
$1/(\Gamma-c\varepsilon)$. It does not show that the learner is slower. \NUM{}

\leanbox{\lean{rate\_floor\_of\_eps\_floor}, \lean{shipped\_floor\_material} in
\lean{LeanPaperA7EtaV1}}

\paragraph{An earlier twelve-run panel} On a centred instance at $M=16$, with $12$ runs at each
$\alpha\in\{0.1,0.05,0.02,0.01\}$ under common random numbers, all $144$ runs returned the correct
answer. The learner's mean was below uniform by $19$--$23\%$ at every level. No cell reached
significance after a Holm correction. \NUM{} The four-cell panel of the main text, with $50$
runs per cell, supersedes it as the measurement of this learner.

\section{S7. The confidence sequence and the martingales}

\paragraph{Ville} A non-negative supermartingale satisfies
$\lambda\Pr(\exists j\le N: M_{j}\ge\lambda)\le\mathbb{E}M_{0}$. Mathlib has Doob's maximal
inequality for submartingales, which is the wrong direction. We prove Ville's inequality from
optional stopping.

\leanbox{\lean{ville\_finite\_horizon}, \lean{confidence\_of\_supermartingale},
\lean{stopping\_error\_le} in \lean{LeanPaperA7ConfidenceV1}}

\paragraph{The closed form at a random past} The learner thresholds
$M_{t}=\prod_{e}\sqrt{a/(a+N_{e})}\exp(S_{e}^{2}/2(a+N_{e}))$. The one-step identity holds at a
fixed past, and under adaptive sampling $S_{e}$ given $N_{e}$ is not Gaussian. A freezing lemma
built from \texttt{condDistrib} licenses the substitution. Integrability is derived by the
product criterion. Crossing the threshold forces $M_{t}\ge1/\delta$, because the penalty
cancels.

\leanbox{\lean{condDistrib\_const\_of\_indepFun}, \lean{condExp\_freeze} in \lean{LeanPaperAFreezeV1};
\lean{martingale\_step}, \lean{joint\_integrable} in \lean{LeanPaperAAdaptMartV1};
\lean{Mval\_eq\_Mstat}, \lean{threshold\_crossing}, \lean{closed\_form\_error\_bound} in
\lean{LeanPaperAStopJoinV1}}

\paragraph{Adaptive queries and the grid prior} An adaptively queried Gaussian stream carries a
martingale mixture. A fixed grid of tilts misses the exponent whenever the true gap is below
half the mesh. A refining mesh loses only the squared miss, since
$u\Delta-\frac12u^{2}=\frac12\Delta^{2}-\frac12(u-\Delta)^{2}$.

\leanbox{\lean{obs\_indep}, \lean{mix\_supermartingale}, \lean{stream\_confidence},
\lean{condExp\_adaptIncr}, \lean{adaptMix\_supermartingale}, \lean{adaptive\_confidence},
\lean{fixed\_grid\_rate\_negative}, \lean{tiltRate\_eq}, \lean{mesh\_grid\_attains\_rate},
\lean{grid\_penalty\_negligible} in \lean{LeanPaperA7StreamV1}}

\section*{Additional proof map: the general repair certificate}
The main paper's certificate proposition uses \lean{exchange}, \lean{frontier\_upper},
\lean{budget} and \lean{weak\_duality} in \lean{LeanPaperASacrificeDualityV1}.
Its assumptions are inequalities for every allocation: a frozen upper certificate, true
alternative loads, and a nonnegative cover of those loads by the frozen margins.
They are not assumed for our example: \lean{cStar\_cert}, \lean{pairLoad\_cert},
\lean{nuStar\_cover} and \lean{certificate\_pair\_satisfiable} construct them, and
\lean{construction\_budget} takes the strict-alternative limit.
The converse uses \lean{repair\_mix} with mixing coefficient
$\lambda=\min\{1,b/s(v)\}$. Therefore its valid budget form is
$R(b)\ge\min\{1,b/s(v)\}J_t(v)$, for $s(v)>0$ and $b\ge0$.
The unclipped expression $bJ_t(v)/s(v)$ requires $b\le s(v)$.
The module's introductory prose omits that restriction, but its actual theorem explicitly
requires the mixing coefficient in $[0,1]$ and is the result used here.
Exact bounds hold for every $r\ge1,M>0$. Statements with absolute constants using
$\Theta(rM)$ or $\Theta(\beta/(rM))$ are restricted to $rM\ge1$; they are not joint limits
as $rM\downarrow0$.

\section{S8. Heteroscedastic families and the repair criterion}

\paragraph{Two families at a fixed design (exact)} With one cheap cell of variance $M$ and all
others unit, two instances show the surrogate mispricing in opposite directions, at the unit
design.
\begin{itemize}
\item Family~1, $A=(0,1,\frac12)$, $B=(0,\frac34,\frac12)$: $J_{\rm full}/J_{\rm fr}\le16/M$.
\item Family~2, $A=(0,1,\frac34)$, $B=(0,\frac12,\frac14)$: $J_{\rm full}/J_{\rm fr}\ge(M+3)/24$.
\end{itemize}
$J_{\rm full}$ is the least of $81$ convex quadratic programs, one per assignment of the
extremum cells. Each is strictly convex and feasible, so it has one minimiser, and the KKT
conditions certify it. Every region of both families is certified in exact rationals at
$M=1,4,16,64,256,1024$. A set of active rows proposed by the floating-point solver, or else found by
exact search over independent subsets, gives multipliers $\mu\ge0$, exact feasibility and
complementary slackness. The exact values are $J_{\rm fr}=\frac1{128}$ and
$J_{\rm full}=\frac{1}{8(M+3)}$ for family~1 at $M\ge16$, and $J_{\rm fr}=\frac{1}{8(M+3)}$ and
$J_{\rm full}=\frac{M+7}{32(4M+3)}$ for family~2 at $M\ge4$. At $M=1024$ the ratios are
$\frac{1024}{65728}=0.0156$ and $64.58$, and both bounds hold exactly at all six points. \CERT{}
Family~2's lower side holds for every $M$ by a direct argument, which is not formalised.
Reversing the ranges needs $a_{3}'-a_{1}'\le b_{j}'-b_{j'}'$ for some pair, and on four unit cells
that costs at least $\frac{(1/4)^{2}}{8}=\frac1{128}$, while $J_{\rm fr}=\frac1{8(M+3)}$. The two directions after the sup over
designs are the theorems of the main text on the eight-cell class
(\lean{high\_separation\_at\_world}, \lean{low\_separation\_at\_world} in
\lean{LeanPaperASeparationV1}). (\texttt{EXP\_paperA\_frozen\_proxy\_exact}, 7/7; floating-point
producer \texttt{EXP\_paperA\_frozen\_proxy}.)

\paragraph{The repair criterion} A design whose surrogate value sits below the barrier to the
nearest switch has $J_{\rm true}=J_{\rm proxy}$. Wherever the criterion fires, the design
inherits the separation. Measured against $\Gamma_{\rm true}$, however, that ceiling is vacuous on
the sweep. At $M=256$ the rational design $\frac{1}{1000}(1,1,1,1,1,1,497,497)$ fires the
criterion at $4.15\%$ of the bracketed true optimum. \CERT{} The value of the best certified
design is not determined.

\leanbox{\lean{barrier\_implies\_equality}, \lean{optima\_equal\_of\_certificate},
\lean{certifiable\_fraction\_le\_sixtyfour\_over\_M}, \lean{design\_certificate\_insufficient} in
\lean{LeanPaperARepairV1}}

\section{S9. The repair law (Theorems ``The repair law'' and ``The frontier is the optimal
rate'', Propositions ``The classes are nested'' and ``Two arms'')}

\paragraph{The classes} A count vector meets the budget in the homogeneous form
$(1-b)\Gamma_f|N|\le\sum_qN_q\mathrm{KL}_q(\theta,\lambda)$ for every frozen alternative, stated in
$\mathbb R_{\ge0}^\infty$, so no division by $|N|$ occurs. Integrating it against the history law
gives $(1-b)\Gamma_f\,\mathbb E|N|\le\mathbb E|N|\cdot\mathrm{info}(\bar w,\lambda)$, by
linearity of the lower integral and finiteness of every count. The infimum over $\lambda$ gives
$J_f(\bar w)\ge(1-b)\Gamma_f$. On the eventual class the pathwise inequality holds up to
$(1-b)\Gamma_ft_0$, since a short history has $|N|<t_0$ and nonnegative information.
\leanbox{\lean{MeetsBudget}, \lean{PrefixClass}, \lean{TerminalClass}, \lean{EventualClass},
\lean{lintegral\_budget\_lhs}, \lean{lintegral\_budget\_rhs}, \lean{terminal\_expected},
\lean{eventual\_expected}, \lean{mean\_escapes} in \lean{LeanPaperAPolicyClassesV1}}

\paragraph{Attainment} The learner is the oracle-tracking learner of Paper A6 with any
measurable allocation-valued oracle of the estimate: warm start, forcing $1/\sqrt{s+1}$,
post-warm quota, first maximiser of $W-N$. Its run is a warm-started tracking run. So the
tracking discrepancy $2(K-1)/(K+k)$ and the count floor $\sqrt t/K-(K-1)$ hold on every path.
The ideal learner's direct bound is re-proved with the oracle's selection step replaced by one
hypothesis: on the confidence event, the oracle's design eventually has true information at
least $\Gamma-\eta/4$. For a constant oracle this is immediate. The rest of the chain is
unchanged:
\begin{itemize}
\item the window average of $(1-\rho_s)(\Gamma-\eta/4)$;
\item the multiplicative tracking factor, the ridge factor and the raw-mean radius;
\item the window failure probability $16/t^2$ at the truth;
\item the layer-cake identity for $\mathbb E\tau$ on $\mathbb N\cup\{\infty\}$.
\end{itemize}
Correctness is Ville's inequality at the stop's own statistic, for every allocation rule.
\leanbox{\lean{o\_count\_floor}, \lean{o\_hpost}, \lean{o\_track}, \lean{o\_radius\_run},
\lean{oracle\_direct}, \lean{o\_goodWin\_compl\_le}, \lean{oracle\_expect\_bound},
\lean{oracle\_correct}, \lean{constO}, \lean{fix\_expect\_bound}, \lean{fix\_upper\_rate} in
\lean{LeanPaperABudgetUpperV1}}

\paragraph{The realised design} With a constant oracle $w$ the average target is at least
$(1-(K+2\sqrt t)/t)w$ cellwise. Multiplicative tracking loses a further $1-2K(K-1)/\sqrt t$. So the
realised design dominates $(1-\varepsilon)w$ after a deterministic time. A game value is
monotone and homogeneous in the design, so the frozen value of the realised design is at least
$(1-\varepsilon)J_f(w)$. Its sacrifice is then at most $s(w)+\varepsilon J_f(w)/\Gamma_f$.
\leanbox{\lean{avgTarget\_const\_ge}, \lean{fix\_design\_floor}, \lean{Jval\_ge\_scale},
\lean{fix\_budget\_eventually} in \lean{LeanPaperABudgetUpperV1}}

\paragraph{Two arms} The specification has unit noise and alternatives
$\{\lambda:\lambda_T\le\lambda_{1-T}\}$. With unit noise the unit-variance coordinates are the
identity, and the energy of a design against an alternative is the two-arm energy. So the
stream game is the closed form $w_1w_2\Delta^2/2$ and the frozen game is $w_1\Delta^2/2$.
\leanbox{\lean{twoSpec}, \lean{jw\_two}, \lean{jf\_two}, \lean{wb}, \lean{two\_arm\_budget\_rate},
\lean{two\_arm\_eventually} in \lean{LeanPaperATwoArmRateV1}}

\paragraph{The rate} $R$ is concave on $[0,\infty)$: take designs within $\varepsilon/2$ of each
frontier value and mix them. So $R$ is continuous at every $b>0$. A design with sacrifice
strictly below $b$ comes within $\varepsilon$ of $R(b)$, through a budget $b'<b$ close to $b$. On
the lower side, the crude bound $R\le J_{\max}$ makes $\mathbb E|N|$ large. Then the excess
$c/\mathbb E|N|$ lies in the continuity ball, and the ideal learner's lower-rate lemma applies at
$R(b)+\eta_0$.
\leanbox{\lean{Rfront}, \lean{R\_le}, \lean{R\_mono}, \lean{R\_concave}, \lean{R\_continuousAt},
\lean{near\_design}, \lean{frontier\_lower\_rate} in \lean{LeanPaperARepairRateV1}}

\paragraph{The frontier inside the law} The budget laws of \lean{LeanPaperARepairLawV1} and
\lean{LeanPaperAPolicyClassesV1} take a bound on $J_{\rm true}$ over every nonnegative vector
within the budget. The frontier bounds it only on the simplex. Since $J_{\rm fr}$ and $J_{\rm true}$
are homogeneous, the design $cv$ has sacrifice $1-cJ_{\rm fr}(v)/\Gamma_f\to-\infty$ and true
information $cJ_{\rm true}(v)\to\infty$, so no finite $R$ meets that hypothesis. The repair: the
expected design is an allocation that puts no mass on the dummy, so the bound is needed only on
the simplex. With an event of probability $p$ at $\theta$ the law reads
$p\log(1/\delta)-\log2\le\mathbb E|N|\,J_{\rm true}(\bar w)$, and on the eventual class
$J_{\rm true}(\bar w)\le R(b+(1-b)t_0/\mathbb E|N|)$ with $R$ the frontier itself. An infimum of
linear forms is concave and bounded by $\sum_q\mathrm{KL}_q(\theta,\lambda_0)$ on designs, which
are the hypotheses of the frontier's concavity and continuity. So a family of $\delta$-correct
policies of the eventual class with one burn-in has
$\liminf\mathbb E|N|/\log(1/\delta)\ge1/R(b)$ on every channel.

\emph{No positivity.} These laws first took $\mathbb E|N|>0$. It is not needed. The expected
informative total is a sum of nonnegative expected counts, so if it vanishes every informative
count vanishes in expectation; the dummy has divergence zero; so the information against every
alternative is zero, and the change of measure gives $p\log(1/\delta)-\log2\le0$, which is the
law with a zero right side. For a $\delta$-correct family, positivity below $\delta<\frac18$ is a
consequence, since there $D_\delta>0$.
\leanbox{\lean{etot\_nonneg}, \lean{info\_zero}, \lean{event\_law0}, \lean{budget\_law\_alloc0},
\lean{eventual\_frontier\_law0}, \lean{frontier\_rate\_lower0} in \lean{LeanPaperARepairNoPosV1};
\lean{etot\_pos\_of\_correct} in \lean{LeanPaperARepairGaussV1}}
\leanbox{\lean{ebar\_sum}, \lean{ebar\_dummy}, \lean{ebar\_alloc}, \lean{event\_law},
\lean{budget\_law\_alloc}, \lean{eventual\_frontier\_law}, \lean{infoAt\_mix},
\lean{Jt\_concave}, \lean{Jt\_le\_sum}, \lean{frontier\_rate\_lower} in
\lean{LeanPaperARepairBridgeV1}}

\paragraph{The law at the stopping time} Meeting the budget at the stop says nothing about a
truncation, so the law for unbounded stopping does not go through truncated budgets. A stopping
policy (kernels for every round, a measurable stop set after every number of rounds, an output
read at the stop) queries the dummy after it stops, and one kernel policy serves every horizon.
At a horizon where at most $\varepsilon$ of the runs at $\theta$ are unstopped, the complement of
the event ``stopped with the answer'' lies in the union of the wrong stopped runs and the
unstopped ones, so the event has probability at least $1-\delta-\varepsilon$. The change of
measure at that horizon bounds $(1-\delta-\varepsilon)\log(1/\delta)-\log2$ by the truncation's
expected information. The truncation's expected counts are at most the whole run's, which are
the supremum over horizons. Taking $\varepsilon$ to zero, with the unstopped mass chosen at
$\varepsilon/(\log(1/\delta)+1)$, gives $D_\delta\le\sum_qN_q\mathrm{KL}_q(\theta,\lambda)$ for the
whole run's counts. The dummy is queried forever after the stop, so its count is infinite; its
divergence is zero, and finiteness is asked of the informative queries only. The infimum over
the alternatives and the budget at the whole run's expected design finish the proof. A correct
finite-horizon policy, stopped at its horizon with the indicator of its answer event, meets every
hypothesis: before the horizon it has not fired, after it the event of the first $N$ rounds keeps
its probability, and its informative counts are constant after $N$ because the query kernel is
the Dirac mass at the dummy. The same zero-total argument removes positivity from the stopped law.
\emph{Satisfiability, budget included.} Embed the stopped fixed-design learner of the optimal-rate
theorem, at its Markov horizon, as a stopping policy. Its whole-run expected counts are those of
the finite-horizon policy, so its expected design and informative total are the same. That policy
lies in the eventual class with one burn-in $t_0$, so its expected design has sacrifice at most
$b+(1-b)t_0/\mathbb E|N|$, and $\mathbb E|N|>0$ below $\delta<\frac18$. Hence every hypothesis of the
stopped law holds with that budget, for every $\delta<\frac18$; on two Gaussian arms for every
$b\in(0,\frac12]$.
\leanbox{\lean{stopAns}, \lean{cntSt}, \lean{ecnt\_le\_St}, \lean{trunc\_event},
\lean{stopped\_info}, \lean{etotSt}, \lean{ebarSt}, \lean{ebarSt\_alloc}, \lean{stopped\_budget\_law},
\lean{finStop}, \lean{finStop\_restr}, \lean{cnt\_mono}, \lean{finStop\_cnt\_after},
\lean{finStop\_cnt\_before}, \lean{finStop\_cntSt}, \lean{finStop\_hyps} in
\lean{LeanPaperARepairStopV1}; \lean{etotSt\_nonneg}, \lean{infoSt\_zero},
\lean{stopped\_budget\_law0}, \lean{finStop\_term}, \lean{finStop\_etotSt}, \lean{finStop\_ebarSt},
\lean{stopped\_law\_satisfiable}, \lean{stopped\_law\_satisfiable\_two} in
\lean{LeanPaperARepairNoPosV1}}

\paragraph{The optimal rate, on one class} The channel of a specification with Gaussian cells is
the Gaussian channel on the cells plus a fresh dummy. Its per-query divergence is Mathlib's
Gaussian $\mathrm{KL}$, $(\theta_e-\lambda_e)^2/(2\sigma_e^2)$, so the kernel information of a
design is the energy of its cell part in unit-variance coordinates, and the kernel game of the
lifted alternatives is the stream game, as infima of the same set. A design with mass $p<1$ on
the dummy is renormalised: its cell part divided by $1-p$ has both games divided by $1-p$, by
homogeneity. The stopped fixed-design learner is a kernel policy through the deterministic
policy of its rule, and its transcript law is the image of the stream measure under the stopped
run. At horizon $n$ its cell counts are the stream counts at some $m\le\min(n,\tau)$, and they
sum to $m$ because every round queries one cell. After the burn-in of the design floor the
realised design meets the budget on every path, which gives the homogeneous inequality for every
frozen alternative at every horizon. The count vector takes values in a countable set, the image
of $\mathbb N^{Q}$, so the event is measurable and the pathwise statement transfers to the
transcript law. Markov's inequality gives a horizon where at most $\delta/|\mathrm{Ans}|$ of the
runs at $\theta$ are unstopped; the test run at level $\delta/|\mathrm{Ans}|$ is then correct at
level $\delta$ at $\theta$ against its alternatives. The horizon depends on $\theta$, like the
design, so the truncated policy is a comparison policy for the optimal cost at $\theta$, not a
learner correct at every world. The untruncated learner stops with a wrong answer with
probability at most $(|\mathrm{Ans}|-1)\delta$ at every world with a unique correct answer
(Ville); its termination is proved at $\theta$ only. The expected informative count at the horizon is at most $\mathbb E\tau$, and the
learner's expected-stopping bound, with $\log(|\mathrm{Ans}|/\delta)=\log(1/\delta)+\log|\mathrm{Ans}|$,
gives the upper rate. On the lower side a correct policy has a positive informative count:
otherwise every informative expected count vanishes, the dummy carries no divergence, and the
change of measure gives $D_\delta\le0$. On two Gaussian arms the frozen optimum $(1,0)$ attains
$\Gamma=\Delta^2/2$, the answer is unique, and the design $(1-b,b)$ within the budget has true
information $b(1-b)\Delta^2/2>0$, so every hypothesis holds.
\leanbox{\lean{gC}, \lean{klq\_gC\_inl}, \lean{klq\_gC\_ne\_top}, \lean{infoAt\_gC}, \lean{Jt\_gC},
\lean{Jval\_smul}, \lean{normalize\_design}, \lean{counts\_meas}, \lean{ae\_map\_counts},
\lean{fixPol}, \lean{hist\_fixPol}, \lean{sum\_Nc}, \lean{real\_counts}, \lean{fix\_eventual},
\lean{fix\_etot\_le}, \lean{etot\_pos\_of\_correct}, \lean{fix\_correct}, \lean{optimal\_lower},
\lean{optimal\_upper}, \lean{optimal\_satisfiable} in \lean{LeanPaperARepairGaussV1}}

\paragraph{Forcing, for unbounded stopping} The truncation of a policy that stops almost surely
at $\theta_C$, at a horizon where at most $\delta$ of its runs are unstopped, is correct at level
$2\delta$. The finite-horizon forcing law applies to it, and a nondecreasing $\varphi$ carries the
bound from the truncation's expected real queries to any bound on the whole run's.
\leanbox{\lean{forcing\_law\_stopped}, \lean{power\_forcing\_law\_stopped} in
\lean{LeanPaperAForcingStopV1}}

\section{S10. The build}

\paragraph{The last full gate} \texttt{full7\_20260925} built $219$ modules and audited $1704$
theorems, with exit status $0$, no \texttt{sorry}, no \texttt{native\_decide}, every axiom closure
within the standard three, and frozen sources byte-identical to the working copy afterwards.
It covers every module cited in the paper and in this supplement except the thirteen listed next.

\paragraph{Smoke gates since} Each module written after \texttt{full7\_20260925} passes a smoke
gate that freezes the module with its full import closure and applies the same acceptance
conditions, with no drift.
\begin{itemize}
\item \lean{LeanPaperASacrificeStoppedV1}: \texttt{smokeR2\_20260925} ($96$ modules, $771$ theorems).
\item \lean{LeanPaperARepairLawV1}: \texttt{smokeUP1\_20260925} ($79$ modules, $762$ theorems).
\item \lean{LeanPaperAPolicyClassesV1}, \lean{LeanPaperABudgetUpperV1}: \texttt{smokeUP2\_20260925}
($107$ modules, $934$ theorems).
\item \lean{LeanPaperATwoArmRateV1}, \lean{LeanPaperARepairRateV1}: \texttt{smokeUP3\_20260925}
($106$ modules, $924$ theorems).
\item \lean{LeanPaperARepairBridgeV1}, with \lean{LeanPaperASacrificeDualityV1} in its closure:
\texttt{smokeUP4\_20260925} ($139$ modules, $1199$ theorems).
\item \lean{LeanPaperARepairStopV1}, \lean{LeanPaperARepairGaussV1}, \lean{LeanPaperAForcingStopV1},
\lean{LeanPaperARepairSlopeV1}: \texttt{smokeUP5\_20260925} ($106$ modules, $835$ theorems).
\item \lean{LeanPaperARepairNoPosV1}: \texttt{smokeUP6\_20260925} ($102$ modules, $825$ theorems).
\end{itemize}
Every citation of the paper and this supplement resolves to a declaration in a module whose
exact source, by hash, passed one of these gates. All the modules above are in the module list of
the next full gate.

\end{document}
